All 412 raw frames were reduced without failures on a laptop (8 cores, 8\,GB of memory), four frames in parallel.
Table~\ref{tab:timing} gives the wall-clock time per frame. Arcs dominate: the 2-d wavelength solution of an arc takes
$\approx48$\,s ($\approx34$\,s since the faster line fits of Section~\ref{sec:visit_timing}), against 4\,s for a lamp flat and 11--16\,s for a science or standard exposure. With improved v3 a
science exposure takes 10\,s, because it no longer searches for the shift of the fringe pattern along the slit in the
2-d frame. Section~\ref{sec:visit_timing} gives the time per visit and per pipeline stage for improved v3, and the
projection to a decade of data.
Disk space matters as much as processor time. A processed 2-d frame is 47\,MB, of which 43\,MB is the uncompressed
table of individual pixels (\texttt{BINNED2D}). A decade of data therefore needs either fewer products written per
frame or compressed tables.

\begin{table}[!htbp]
\centering
\caption{Processing time per frame (wall-clock, four frames in parallel on an 8-core laptop).}
\label{tab:timing}
\begin{tabular}{lccc}
\toprule
Frame type & $N$ & median time & total, 4 workers \\
 & (frames) & (s) & (min) \\
\midrule
Arc & 122 & 48 & 26 \\
Lamp flat & 133 & 4 & 3 \\
Standard (Feige\,110) & 39 & 16 & 3 \\
Science (Fairall~9, EG\,21) & 118 & 11 & 6 \\
Science, improved v3 & 118 & 10 & 6 \\
Standard, improved v3 (two passes) & 39 & $2\times12$ & 4 \\
\bottomrule
\end{tabular}
\end{table}
